\section{Security Analysis}
\label{sec:security-analysis}

\subsection{Authorization Observability}

For execution context $u$, let $\mu(u)$ be its committed effects and let
$\gamma(u)$ contain policy-relevant metadata such as effect provenance.  The
ideal commit view is $\omega(u)=(\mu(u),\gamma(u))$.  Each policy context
$P\in\mathcal P$ induces an ideal decision $I_P(\omega(u))$.  Our main
experiment instantiates $P$ with ACLs, capabilities, and delegations; a finite
stress relation uses effect-set containment. The checked policy snapshot also
governs the commit. A monitor
sees $\rho(u)$ rather than $\omega(u)$; a descriptor instantiates
$\rho_d(u)=A_{d_t}(u)$.

Representation induces the observational equivalence relation
$u\sim_\rho v$ exactly when $\rho(u)=\rho(v)$.  Every downstream authorizer
whose call-dependent input is $\rho$ is constant within each such class.  The
interface therefore upper-bounds the resolution of every downstream policy,
independently of the policy language's sophistication.

\begin{definition}[Authorization sufficiency]
\label{def:authorization-equivalence}
A representation $\rho$ is authorization-sufficient on domain $D$ for policy
family $\mathcal P$ when
\[
 \rho(u)=\rho(v) \Longrightarrow
 I_P(\omega(u))=I_P(\omega(v))
\]
for every $u,v\in D$ and $P\in\mathcal P$.
\end{definition}

The definition applies to every representation. It permits distinctions ignored
by $\mathcal P$ to remain merged and requires every policy-distinct pair to be
separated.

\begin{theorem}[Authorization-observability collision]
\label{thm:representation-insufficiency}
If $\rho$ is not authorization-sufficient, no deterministic monitor whose only
call-dependent input is $\rho(u)$ can, even when given $P$, both allow every
authorized context and reject every unauthorized context.
\end{theorem}

\paragraph{Proof sketch.}
Insufficiency gives $u,v,P$ with equal representations and opposite ideal
decisions.  The monitor must return the same decision for both.  Allowing
admits the unauthorized member; denying or abstaining withholds the authorized
member.

For fixed $P$, if a representation cell contains both an authorized and an
unauthorized row, every sound monitor over that view must deny or abstain on
the authorized row, yielding at least one false denial or abstention. This bound applies to tool-name, whole-call, and atom
views alike.

This information bound applies equally to ACL, capability, provenance, and
model-based policy consumers. Agent harnesses make it directly testable:
schemas, defaults, aliases, and pre-state determine what a generated call
commits, while authorization acts on the interface exposed before commitment.
Validation exposes disagreements between these views as
implementation-grounded counterexamples.

\subsection{Executable Interface Falsification}

An intervention $v=do_J(u)$ changes selected call fields, trusted provenance,
or pre-state while holding the tool implementation fixed.  Source execution
observes committed state changes; the trusted harness records policy-relevant
provenance and control metadata.  Together they instantiate $\omega$ on the
tested pair.

\begin{proposition}[Counterfactual witness]
\label{prop:mediation-gap}
If $\rho(u)=\rho(v)$ but some $P\in\mathcal P$ satisfies
$I_P(\omega(u))\ne I_P(\omega(v))$, then $\rho$ is insufficient on every domain
containing $u$ and $v$.
\end{proposition}

The proposition turns counterfactual testing into a falsification procedure.
An output-only change with unchanged $\omega$ provides no witness. A trusted
provenance change may provide one when policy distinguishes that provenance;
sensitivity to a semantic-invariant placebo instead exposes overpartition.

An effect witness and a policy-separating witness form a failure certificate
when the candidate representation merges the pair. Source execution, policy
separation, and representation equality supply its three predicates.

For a finite domain, counterexample-guided refinement repeatedly splits a
representation cell using a witnessed separating field or qualifier.  Call a
cell \emph{policy-mixed} when it contains $u,v$ for which at least one
$P\in\mathcal P$ gives different ideal decisions.

\begin{proposition}[Finite refinement termination]
\label{thm:finite-adequacy}
On finite $D$, a refinement procedure that never merges cells can perform at
most $|D|-1$ nonempty cell splits.  If each step separates a witnessed
authorization collision and the procedure stops only when no policy-mixed cell
remains, the resulting representation is authorization-sufficient on
$(D,\mathcal P)$.
\end{proposition}

\paragraph{Proof sketch.}
Every split strictly increases the number of nonempty cells, which is bounded
by $|D|$.  At termination, equal representations imply membership in the same
nonmixed cell.  By the definition of policy-mixed, their ideal decisions agree
for every $P\in\mathcal P$.

Generated interventions define $D$, and unresolved interventions remain in the
validation record. Each accepted descriptor carries the $(D,\mathcal P)$ pair
on which the adequacy result applies.

\subsection{Conditional Consumer Composition}

The representation results become a commit guarantee only when the descriptor,
policy, and enforcement path satisfy separate obligations.  Let $J_P$ be the
ideal policy predicate expressed over a descriptor's atom view.

\begin{proposition}[Conditional interface consumption]
\label{thm:conditional-atom-mediation}
Assume four obligations for every reachable effectful call $u$.  First,
descriptor fidelity gives $J_P(A_{d_t}(u))=I_P(\omega(u))$ for every admissible
$P$.  Second, the authorizer returns \Allow{} only when
$J_P(A_{d_t}(u))=1$.  Third, every effectful call is mediated before commit.
Finally, the executor commits exactly the checked call under the same policy
snapshot.  Then every allowed committed call satisfies $I_P(\omega(u))=1$.
\end{proposition}

\paragraph{Proof sketch.}
Consider an allowed committed call $u$.  Complete mediation supplies a check
for $u$, and authorizer soundness gives $J_P(A_{d_t}(u))=1$.  Descriptor
fidelity therefore gives $I_P(\omega(u))=1$.  Call and policy-snapshot
check--use consistency make this conclusion apply at commitment.

Counterfactual tests probe descriptor fidelity. The explicit-authority
experiment consumes the resulting interface through one ACL, capability, and
delegation engine; AgentDojo instantiates a provenance-origin policy over
registered fields.

\subsection{Integration Instance: Provenance-Origin Confinement}

For this integration policy, let $F(u)$ and $E(u)$ be the registered field values
and effect kinds of call $u$.  Let $U_V(h),U_E(h)$ denote values and requests
from provenance-marked untrusted evidence, and let $G_V(q),G_E(q)$ denote those
grounded in the authenticated task.  The monitor allows only when
\[
  F(u)\cap U_V(h)\subseteq G_V(q)
  \quad\text{and}\quad
  E(u)\cap U_E(h)\subseteq G_E(q).
\]

\begin{corollary}[Conditional provenance-origin confinement]
\label{thm:single-commit}
If the descriptor, provenance adapter, and grounding relation are sound for the
registered fields and effect kinds, and every committed call is exactly the
call checked, then an allowed call contains no registered effect or value
introduced solely by untrusted evidence outside the authenticated task.
\end{corollary}

\paragraph{Proof sketch.}
Any violating value lies in $F(u)\cap U_V(h)$ but outside $G_V(q)$, contradicting
the first condition; the same argument over $E(u)$ contradicts the second for
an effect-only violation.  Exact mediation transfers the check to commitment.

This corollary characterizes the integration policy.  The explicit-authority
experiment evaluates ACL, capability, and delegation decisions separately;
runtime audits here test mediation and check--use consistency.
Appendix~\ref{app:formal-proofs} gives the complete definitions and proof
details.
